为突出红外干涉测厚的主要物理机制，并保持各问简化条件一致，作出以下假设：

\begin{enumerate}[label=\textbf{假设\arabic*:},leftmargin=4.5em,itemsep=0.35ex]
\item \textbf{理想分层结构。} 外延层在光斑范围内厚度和材料参数近似均匀，空气--外延层及外延层--衬底界面光滑、平行，材料可视为各向同性层状介质，因此采用一维平行分层传播模型。

\item \textbf{测量状态稳定。} 同一晶圆在$10^\circ$和$15^\circ$测量期间，样品温度、波数标定和入射角控制保持稳定，环境变化造成的折射率漂移和热膨胀相对厚度反演尺度可忽略。

\item \textbf{非偏振独立传播。} 入射光视为非偏振光，界面不产生显著偏振耦合，$s$、$p$两种偏振分别按Fresnel关系计算，非偏振反射率取两者反射强度的算术平均。

\item \textbf{波段内光学模型适用。} SiC在$2500$--$3300\,\mathrm{cm}^{-1}$弱吸收区采用实折射率近似；硅在$1500$--$3500\,\mathrm{cm}^{-1}$内保留复折射率与传播衰减，并选取满足$0<|P|\le1$的被动介质支路。

\item \textbf{同晶圆厚度一致。} 同一晶圆在两个入射角下对应同一真实外延层厚度，测量区域间的横向厚度差可忽略，因此两角度反演差异主要反映模型近似、噪声和角度条件差异。

\item \textbf{窄波段参数稳定。} 在各自拟合波段内，衬底等效折射率近似为常数；硅束缚电子背景参数和有效质量采用固定文献值，自由载流子变化主要由Drude项描述，以避免在有限光谱内同时引入过多近共线参数。
\end{enumerate}